\ficha{Gremaud (1997), Controvérsias monetárias no Brasil agroexportador}{gremaud1997}

\begin{identificacao}
GREMAUD, Amaury. \emph{As controvérsias monetárias no Brasil
agro-exportador}. Ribeirão Preto: FEA-USP-RP, 1997. 31 p. Texto não
publicado em periódico.

Artigo de história do pensamento econômico, 31 páginas, em português. Lido
integralmente. O PDF consultado não traz veículo, volume nem data de
publicação, e informa em nota de rodapé que é ``uma versão modificada do
capítulo 2 da tese de doutorado do autor sem as adequações às normas da
ABNT'' (p. 1). A data de 1997 vem da semente bibliográfica e não do documento.
\end{identificacao}

\bloco{Problema e tese}
O artigo pergunta como evoluiu o pensamento monetário brasileiro do Império à
Primeira República e como se organizaram os campos em disputa. A resposta é
uma taxonomia com um eixo principal e três subordinados. O eixo principal é o
grau de vinculação da moeda nacional à moeda internacional, que se fazia por
dois elementos, o lastro sobre o qual a moeda era emitida e sua
conversibilidade (p. 1). Os subordinados são a responsabilidade pela emissão,
entre Tesouro, banco monopolista e pluralidade de bancos, e o nível adequado da
taxa cambial (p. 1). A tese é que as posições se moveram: metalistas e
papelistas radicalizaram, inverteram argumentos e, ao final, convergiram para a
defesa de organismos que controlassem o mercado cambial e desenvolvessem o
mercado financeiro nacional (p. 2).

\chave{A controvérsia central foi acerca do grau de vinculação da moeda
nacional à moeda internacional.}{p. 1}

\bloco{Argumento}
\begin{enumerate}[leftmargin=*,nosep]
\item As duas posições de partida. Os metalistas defendem estabilidade cambial
e padrão metálico, e por isso querem unidade emissora, a princípio sob controle
do governo, e só depois da recuperação do caráter metálico admitem a emissão
conversível por um banco único, regulada pela entrada de metais (p. 1). Os
papelistas contrapõem a estabilidade cambial às necessidades da circulação
interna e defendem, no início, emissões conversíveis mas não integralmente
lastreadas, com oferta seguindo a demanda do mercado e não as reservas do país,
por meio da pluralidade bancária com supervisão pública (p. 1-2).

\item Os núcleos teóricos. Para o metalista, o valor da moeda depende de sua
quantidade, o ajuste é automático pela fusão e pela cunhagem, o câmbio é a
variável-chave e reflexo do valor da moeda, e o receituário é a contenção
monetária para que os metais reocupem o espaço do papel, segundo a lei de
Gresham (p. 4-6). Para o papelista, a moeda não tem valor intrínseco, seu valor
vem da função de intermediar trocas, e a referência básica é a taxa de juros, e
não a de câmbio (p. 6-7, n. 3). A contracrítica papelista é que o câmbio é
determinado pela balança comercial e pelo balanço de pagamentos, o que o
desqualifica como indicador de política monetária (p. 7).

\item O deslocamento papelista. A corrente se radicaliza, abandona a moeda
conversível e propõe moeda puramente fiduciária lastreada em títulos internos.
Com isso perde o mecanismo automático de controle, e o resultado é o abandono
da pluralidade emissora em favor de um grande organismo emissor, com a função
adicional de proteger o restante do sistema monetário (p. 2). Souza Franco, em
1848, ainda propunha bancos regionais com fundo de dois terços em apólices e um
terço em ouro, notas conversíveis e emissão limitada ao tamanho do fundo (p. 9),
e recusava o banco central por argumento político, pois criaria no Estado
``corpo excessivamente forte e incompatível com o sistema representativo''
(p. 8). Mauá, em 1878, e Amaro Cavalcanti, em 1893, defendem a moeda
inconversível e sem lastro metálico, afastando-se do princípio bancário inglês
(p. 13). Cavalcanti acrescenta o argumento da dependência: o metal amoedado está
monopolizado por três ou quatro nações, que ditam a lei às demais (p. 14-15).

\item O deslocamento metalista. As dificuldades práticas dividem o campo entre
moderados, que aceitam a meta metálica no longo prazo e a taxas menos
valorizadas que o par do Segundo Império, e radicais, que querem o retorno ao
padrão anterior com deflação vigorosa (p. 2). No fim da Monarquia, Ouro Preto,
o Visconde do Cruzeiro e Lafayette Rodrigues, que Gremaud chama de metalistas
realistas, propõem pluralidade e moeda fiduciária inconversível sem defender
suas vantagens, apenas apontando os obstáculos ao curso metálico (p. 17-18). Sob
pressão dos ortodoxos e com o câmbio alcançando a paridade fixada em 1846, o
projeto foi alterado para incorporar emissão conversível lastreada em metal
(p. 18).

\item Rui Barbosa e a inversão. Sua crítica de 1890 recai justamente sobre a
introdução da conversibilidade e do lastro metálico, porque a nota bancária
ficava ``subordinado ao câmbio de 27'' e teria de trocar ouro ao par notas
cotadas abaixo dele (p. 19). Para os metalistas era a inconversibilidade que
instabilizava o câmbio; para Rui era a instabilidade cambial, de causas não
monetárias, que inviabilizava a moeda metálica (p. 20). Seu barômetro do
excesso de emissão é a taxa de juros, não o câmbio (p. 21). A pluralidade foi
compromisso com o federalismo e não doutrina, e sua preferência é a
centralização, com apoio em Wagner: um grande banco resiste a crises, age de
modo contracíclico e atua no mercado cambial contra a especulação dos bancos
internacionais (p. 22). O resultado é uma política monetária discricionária
sobre base inconversível (p. 22-23).

\item Murtinho e a equação. A partir de 1898 o câmbio volta a ser a
variável-chave e sua valorização, o objetivo básico (p. 24). Murtinho admite que
a emissão possa ser agente de progresso se aplicada produtivamente, e separa o
papel conversível, cuja conversão cria responsabilidade no emissor, do curso
forçado, em que o interesse não tem freio (p. 24-26). A taxa de câmbio é o
quociente entre o papel-moeda em circulação, em moeda local, e as exportações,
em moeda estrangeira (p. 26), que Gremaud lê como espécie de teoria quantitativa
em que as exportações fazem as vezes de renda e o câmbio, de preços (p. 26,
n. 10). Elevar o valor da exportação seria preferível e é lento demais, o que
justifica o resgate e a queima do papel (p. 26-27).

\item A crítica contemporânea e o desfecho. Vieira Souto ataca Murtinho em
artigos no jornal O Correio da Manhã, entre novembro e dezembro de 1901,
em dois pontos, a equação do câmbio e os efeitos depressivos da retirada do
meio circulante, e atribui a valorização em curso ao fluxo de capital, não ao
resgate (p. 27). Depois de Murtinho, o debate se organiza em dois eixos, e a
pergunta sobre qual taxa perseguir separa Leopoldo Bulhões, que queria o padrão
de 1846, de David Campista, para quem aquele patamar era de difícil obtenção e
mesmo equivocado (p. 28). No segundo eixo, carteira de emissão do Banco do
Brasil, carteira de redesconto e transformação em Banco Central, a objeção
principal não é ao poder emissor de um grande banco, mas à inconversibilidade
das notas (p. 29). Antonio Carlos exige lastro de 100\% em ouro e propõe, em
lugar do redesconto, reserva legal obrigatória de 10\% dos depósitos no Banco do
Brasil (p. 29-30); Inglez de Souza admite lastro de 40\% (p. 29-30). Do pós-guerra
resta menos defesa da moeda plenamente metálica e mais atenção à elasticidade e
à resistência a crises, com o grande banco controlador do câmbio ganhando
importância nos dois lados (p. 30).
\end{enumerate}

\chave{Aqui temos a inversão das idéias metalistas.}{p. 20}

\bloco{Conceitos e definições operacionais}
Vinculação à moeda internacional, decomposta em lastro e conversibilidade
(p. 1). É a definição que separa quatro perguntas que costumam andar juntas:
sobre o que se emite, se a nota se troca por metal, quem emite e a que taxa.
Metalismo e papelismo resultam da combinação dessas respostas e não da adesão a
uma escola, e o autor registra que não há homogeneidade dentro de cada campo
(p. 4). Metalistas realistas é o rótulo do autor para quem mantém a meta
metálica e a declara inatingível no presente (p. 18). Valor potencial e valor
real são a distinção com que Murtinho avalia a emissão pelo emprego dado ao
papel (p. 24-25). Câmbio real e câmbio nominal, em Bulhões em 1891 e em
Bernardino de Campos em 1898, separam a parte do câmbio limitada pelo custo de
remessa de metal daquela que acompanha sem limite a depreciação da moeda
(p. 23-24). O autor não define padrão-ouro, credibilidade, regra ou
\emph{currency board}, e não usa esse vocabulário.

\bloco{Evidência e método}
Leitura de fontes impressas de época, sem séries, modelos ou contrafactual.
As fontes são discursos parlamentares, relatórios de ministros da Fazenda,
panfletos e livros de doutrina, entre eles Souza Franco (1848), Mauá (1878),
Cavalcanti (1893), Rui Barbosa (1893), Andrada (1923) e Inglez de Souza (1924),
além de artigos de jornal de Vieira Souto (p. 27). O recorte vai do Primeiro
Reinado ao pós-Primeira Guerra. Boa parte das citações do Império chega por
compilação, com \emph{apud} em Andrada (1923) e em Cavalcanti (1983), e a
leitura da transição republicana se apoia em Franco (1987), citado em bloco
(p. 18-19).

\bloco{Agentes e vertentes}
Metalistas: Visconde de Itaboraí, Torres Homem, Francisco Belisário Soares de
Sousa, Joaquim Murtinho, Leopoldo de Bulhões, Bernardino de Campos, Antonio
Carlos, Calógeras, Ortigão, Inglez de Souza. Papelistas: Souza Franco, Mauá,
Amaro Cavalcanti, Rui Barbosa, Vieira Souto, Pires do Rio. Metalistas
realistas: Visconde de Ouro Preto, Visconde do Cruzeiro, Lafayette Rodrigues.
Sem campo atribuído, na discussão do banco central: David Campista, Carlos de
Campos, Cincinato Braga, Whitaker, Francisco Sá, Homero Batista (p. 28-29,
n. 13). Referências estrangeiras mobilizadas pelos próprios agentes: Gresham,
Wagner, o princípio bancário inglês.

\bloco{Críticas e limites}
O texto não trata da Caixa de Conversão e não a nomeia. Entre a crítica de
Vieira Souto de 1901 e a seção final há um salto, e o período de 1906 a 1914
aparece só por duas linhas sobre a taxa almejada, com Bulhões contra Campista
(p. 28). A seção final também não data as posições: autores que publicam em
1923 e 1924 são postos ao lado de propostas presumivelmente anteriores, o que
torna a cronologia imprecisa. A atribuição de campo é do autor e não vem
acompanhada de critério explícito, tanto que ele admite a heterogeneidade
interna de cada corrente (p. 4). Não há evidência quantitativa, e a
correspondência entre doutrina e política efetivamente adotada é afirmada, não
testada. A dependência de Andrada (1923) e Franco (1987) faz o Império ser lido
por compiladores com agenda própria. O texto é versão de capítulo de tese sem
adequação a normas, com referências incompletas e ao menos uma troca de termo
que inverte o sentido de uma frase sobre a emissão de curso forçado (p. 25-26).

\chave{os quais se faziam sobre dois eixos: o fortalecimento de um grande banco
no intuito de assumir a condição de um quase Banco Central e na instituição de
um regime de câmbio fixo}{p. 28}

\begin{dialogo}
\item Vieira Souto no Correio da Manhã. Gremaud registra artigos entre novembro
e dezembro de 1901 no jornal (p. 27). O Correio da Manhã é um dos quatro
diários do corpus. Procurar, nas páginas do próprio jornal e depois em 1906,
o vocabulário do ``câmbio real'' contra o ``câmbio nominal'', a fórmula da
``relação dos créditos e débitos recíprocos'' e a acusação de que o resgate
enriquece o Tesouro e empobrece ``a lavoura, indústria e comércio'' (p. 27).
Testa se a linha do jornal sobre a Caixa tem genealogia doutrinária na sua
própria página de 1901, ou se é posição formada no calor de 1906.

\item O par de 1846 como marco. Rui Barbosa atacou a conversibilidade porque a
nota ficava subordinada ao ``câmbio de 27'' (p. 19), e quinze anos depois
Bulhões ainda perseguia o padrão de 1846 enquanto Campista o julgava
inatingível e equivocado (p. 28). Procurar nos jornais e nos Documentos
Parlamentares quem ainda invoca o par legal em 1906 e quem o trata como ficção
jurídica. Testa a hipótese do capítulo de que entre 1890 e 1906 a divergência
migrou de ``ouro ou papel'' para ``a que nível fixar'', o que retira da
clivagem metalismo e papelismo a função de organizar o debate da Caixa.

\item Hierarquia das variáveis. O barômetro de Rui é a taxa de juros, não o
câmbio (p. 21), e a referência básica papelista é a mesma (p. 7, n. 3).
Verificar se algum defensor da Caixa entre 1906 e 1914 argumenta pela taxa de
juros, pela escassez de numerário ou pelas necessidades da praça, em lugar do
nível da taxa em pence. Se o argumento de juros estiver ausente em Campista,
Bulhões, Serzedello Corrêa, Alcindo Guanabara e Barbosa Lima, o papelismo
sobreviveu como rótulo e não como quadro analítico.

\item A objeção é ao lastro, não a quem emite. No segundo eixo de Gremaud, a
crítica às carteiras do Banco do Brasil recai sobre a inconversibilidade das
notas, e não sobre a concentração do poder emissor em um grande banco (p. 29),
com Antonio Carlos exigindo 100\% de ouro e Inglez de Souza aceitando 40\%
(p. 29-30). Procurar a mesma gramática na controvérsia sobre a competência da
Caixa frente ao Tesouro e ao Banco do Brasil, sobre o limite de emissão e sobre
o fundo de garantia. Testa se a Caixa foi recebida como dispositivo de lastro
integral na margem, portanto do lado conversível da disputa, ou como embrião de
banco central, e permite datar quando a objeção muda de objeto.

\item O argumento da dependência. Cavalcanti sustenta que três ou quatro nações
monopolizam o metal e ditam a lei às demais (p. 14-15). Procurar, no campo
expansionista de 1906 a 1914, crítica ao ouro formulada em termos de autonomia
nacional diante de Londres e dos Rothschild, distinta do argumento de interesse
setorial do café. A ausência desse vocabulário indicaria que a linha de
Cavalcanti não sobreviveu, e que a defesa da emissão em 1906 se faz por
interesse declarado e não por soberania monetária.
\end{dialogo}

\usonocapitulo{Parte III, como texto que fornece a taxonomia e a cronologia do
movimento das posições até a véspera da Caixa. Serve a três afirmações.
Primeira, que a controvérsia se decompõe em quatro perguntas separáveis, lastro,
conversibilidade, quem emite e nível da taxa (p. 1), o que autoriza descrever
1906 sem forçar o par metalismo e papelismo. Segunda, que as posições se
moveram e chegaram a se inverter, com Rui trocando a ordem causal entre moeda e
câmbio (p. 20), com os metalistas realistas admitindo a moeda inconversível
(p. 17-18) e com os dois campos convergindo para um grande banco controlador do
câmbio (p. 2, p. 30), convergência que é o antecedente imediato da Caixa.
Terceira, que a pergunta sobre a taxa almejada, Bulhões contra Campista sobre o
padrão de 1846 (p. 28), é a mesma que o debate de 1906 responde ao fixar 15
pence, e que os dois nomes já estão postos em campos distintos antes do Convênio
de Taubaté. O texto não trata da Caixa de Conversão, não a menciona e não cobre
a operação do período, de modo que nenhuma afirmação sobre 1906 a 1914 pode se
apoiar nele além das duas linhas sobre a taxa almejada.}
